\documentclass{article}

\usepackage{booktabs}
\usepackage{tabularx}
\usepackage{parskip}
\usepackage{hyperref}

\setlength{\parindent}{0pt}

\title{Development Plan\\\progname}

\author{\authname}

\date{}

\input{../Comments.text}
\input{../Common.text}

\begin{document}

\maketitle

\newpage{}

\begin{table}[hp]
\caption{Revision History} \label{TblRevisionHistory}
\begin{tabularx}{\textwidth}{llX}
\toprule
\textbf{Date} & \textbf{Developer(s)} & \textbf{Change}\\
\midrule
2026/09/25 & Mohammad Bilal & Completed sections 1-4, Team Charter, \& reflection\\
2026/09/25 & Jake Finlay & Completed Section 10: Proof of Concept Demonstration Plan \\
\bottomrule
\end{tabularx}
\end{table}

\newpage{}

\section{Introduction}

This document outlines the team's development plan for the project.
Sections 2--4 cover confidentiality, intellectual property, and copyright,
while Sections 5--10 describe team processes, workflow, scheduling, and the proof of concept.
Sections 11--14 cover the expected technologies, use of AI and LLMs, and coding standards.
The Team Charter outlines the team's expectations for attendance, accountability, teamwork, and decision making.

\section{Confidential Information?}

There is no confidential industry information involved in this project.
However, the team may interact with patients, parents, clinicians, and other
stakeholders during the development and evaluation of the application.
Any personal, medical, or identifying information shared during these interactions
will be treated as confidential and will not be included in the project repository
or documentation in an identifiable form. The team will follow the privacy and
confidentiality requirements provided by the project supervisors and McMaster
University when collecting, handling, or discussing stakeholder information.

\section{IP to Protect}

The intellectual property created for this project will solely consist of the source code,
documentation, and application content developed by the team. These materials will be
protected under the copyright license described in Section~\ref{sec:copyright-license}.
No additional intellectual property protection, such as a trademark agreement, is
currently anticipated for the project.

\section{Copyright License}
\label{sec:copyright-license}

The copyright license adopted by the team for this project is the MIT License.
This license allows others to use, copy, modify, merge, publish, distribute, sublicense,
and sell copies of the software, provided that the original copyright notice and permission
notice are included. The \href{https://github.com/JDBeach1503/hindsight/blob/main/LICENSE}{copyright license}
for the project can be found in the team repository.

\section{Team Meeting Plan}

\wss{This section is for meetings between the team members, and meetings between 
the team and their supervisor/advisor (as appropriate).}

\wss{How often will the team meet? where? when?}

\wss{If the meeting is a physical location (not virtual), out of an abundance of
caution for safety reasons you shouldn't put the location online}

\wss{To help schedule meetings, you can use the lecture/tutorial time when we do 
not have anything else scheduled.}

\wss{How often will you meet with your supervisor/advisor?  when?  where?}

\wss{Will meetings be virtual?  At least some meetings should likely be
in-person.}

\wss{How will the meetings be structured?  There should be a chair for all meetings.  
There should be an agenda for all meetings.}

\section{Team Communication Plan}

\wss{Issues on GitHub should be part of your communication plan.}

\section{Team Member Roles}

\wss{You should identify the types of roles you anticipate, like notetaker,
leader, meeting chair, reviewer.  Assigning specific people to those roles is
not necessary at this stage.  In a student team the role of the individuals will
likely change throughout the year.}

\section{Workflow Plan}

\begin{itemize}
	\item How will you be using git, including branches, pull request, etc.?
	\item How will you be managing issues, including template issues, issue
	classification, etc.?
  \item Use of CI/CD
\end{itemize}

\section{Project Decomposition and Scheduling}

\begin{itemize}
  \item How will you be using GitHub projects?
  \item Include a link to your GitHub project
\end{itemize}

\wss{How will the project be scheduled?  This is the big picture schedule, not
details. You will need to reproduce information that is in the course outline
for deadlines.}

\section{Proof of Concept Demonstration Plan}

The proof of concept (PoC) demonstration will focus on one complete ``vertical
slice'' of the application instead of covering many procedures at once. The team
will demonstrate the following on both a physical iOS device and a physical
Android device:

\begin{enumerate}
  \item A child-friendly walkthrough of one common procedure (receiving eye
  drops), shown through a guide character with narration, simple language, and
  illustrations or animations.
  \item An interactive vision testing practice activity, where the child matches
  picture-based symbols similar to those used on paediatric vision charts.
  \item Navigation between these activities that a child could complete with
  little or no help from a parent.
\end{enumerate}

\noindent This slice was chosen because it exercises every major risk to the project.
The accuracy of clinical content, whether young children can actually use and
engage with the application, and whether rich media can be delivered reliably
on both mobile platforms.

\subsection{Risk 1: Clinical Accuracy and Age-Appropriateness of Content}

The value of the application depends on children getting an accurate preview of
what will happen at their appointment. If a procedure is shown incorrectly (for
example, steps in the wrong order, leaving out that eye drops may sting, or
equipment that looks different from what is in the clinic), the child's
expectations will not match reality. This could increase anxiety, which is the
opposite of the project's goal. No team member has clinical training, so there
is a risk that the team cannot produce content that is both clinically accurate
and understandable for young children. There is also a risk that getting
content reviewed by a clinician takes too long to cover the full set of planned
procedures.

\subsubsection*{What the PoC will demonstrate:}
\begin{itemize}
  \item The eye drop walkthrough (script and visuals) will be reviewed by a
  paediatric ophthalmology clinician before the demonstration. The team will
  record the corrections requested and how many review rounds were needed.
  \item The narration will use short sentences and familiar words, and will be
  honest about sensations (e.g., drops may feel cold or sting briefly) while
  staying reassuring.
  \item The team will measure how long one procedure takes to go from draft to
  clinician-approved. This will be used to estimate how many procedures are
  realistic for the final product.
\end{itemize}

\subsubsection*{Success criteria:}
The clinician confirms the walkthrough accurately
reflects what a child experiences, with no outstanding corrections, and the
review turnaround is short enough to cover the remaining planned procedures
(e.g., vision testing and slit-lamp examination).

\subsection{Risk 2: Child Engagement and Usability}

The target users are young children, many of whom cannot read yet, have short
attention spans, and are still developing fine motor skills. If they cannot
navigate the application on their own or lose interest partway through, they
will not benefit from the preparation. Designing for children is different from
designing for adults and requires audio-first instructions, large touch targets,
and interactions that are forgiving of mistakes. A related risk is that testing
directly with children may require research ethics approval, which may not be
in place by the PoC date.

\subsubsection*{What the PoC will demonstrate:}
\begin{itemize}
  \item All instructions in the vertical slice are narrated, so no step
  depends on the user reading text.
  \item Touch targets are sized for small fingers, and incorrect taps in the
  vision practice activity lead to encouragement instead of failure.
  \item Informal walkthroughs will be run with people unfamiliar with the
  application (e.g., parents or peers outside the team). The team will record
  whether they complete the walkthrough and activity without verbal help, how
  long it takes, and where they get stuck. If ethics approval is in place, a
  supervised session with a child will also be run.
  \item The team will confirm what approval is required for testing with
  children and a realistic timeline for obtaining it.
\end{itemize}

\subsubsection*{Success criteria:}
Testers complete both activities without verbal
help, the full slice takes about five minutes or less so that it fits a young
child's attention span, and the team has a clear path to testing with children
before the final evaluation.

\subsection{Risk 3: Cross-Platform Delivery of Rich Media}

The application must run on iOS and Android, on both phones and tablets, while
playing animations, audio narration, and interactive touch content. Families may
use older or lower-end devices, and the application may be opened in hospital
waiting rooms where connectivity is poor. As more procedures are added, media
files could make the application large to download. With a small team, keeping
one codebase working well on both platforms is a technical risk.

\subsubsection*{What the PoC will demonstrate:}
\begin{itemize}
  \item The same build of the vertical slice running on at least one iOS device
  and one Android device, including at least one older or lower-spec device.
  \item Narration and animations playing smoothly and in sync on each device.
  \item The slice working with no network connection (airplane mode) after
  installation.
  \item The install size of the slice, used to estimate the total size of the
  application once all planned procedures are included.
\end{itemize}

\subsubsection*{Success criteria:}
The slice behaves the same on both platforms with no
noticeable lag or audio/visual desynchronization, works fully offline, and the
estimated full application size is reasonable for a family to download.

\subsection{Consequences of an Unsuccessful PoC}

If the PoC does not sufficiently address these risks, the consequences for the
project are as follows.

\subsubsection*{If clinical accuracy (Risk 1) is not resolved:}
\begin{itemize}
  \item The team cannot claim the content is trustworthy. Inaccurate content
  could mislead children and increase their anxiety, and clinicians would be
  unlikely to recommend the application to families.
  \item If clinician review is too slow, fewer procedures can be covered in the
  final product (e.g., only eye drops and vision testing, without the slit-lamp
  examination), leaving parts of the appointment unprepared for.
\end{itemize}

\subsubsection*{If child engagement and usability (Risk 2) are not resolved:}
\begin{itemize}
  \item If children cannot use the application on their own, it becomes a
  parent-led tool. This would require rethinking the interaction design and may
  make the application less engaging for the child.
  \item If approval to test with children cannot be obtained in time, the final
  evaluation would have to rely on adult testers and clinician judgement. This
  would weaken any claims that the application actually helps children feel
  more prepared.
\end{itemize}

\subsubsection*{If cross-platform delivery (Risk 3) is not resolved:}
\begin{itemize}
  \item Supporting only one platform would exclude families who use the other,
  reducing how accessible the application is.
  \item Poor performance on older devices, or a very large download, would
  create barriers for some families. The team may need to reduce the richness
  of the media (e.g., static illustrations instead of animations), which could
  make the application less engaging for children.
\end{itemize}

\noindent Whether or not the PoC succeeds, its results will be used to set the number of
procedures covered, the interaction design, and the media approach.

\section{Expected Technology}

\wss{What programming language or languages do you expect to use?  What external
libraries?  What frameworks?  What technologies.  Are there major components of
the implementation that you expect you will implement, despite the existence of
libraries that provide the required functionality.  For projects with machine
learning, will you use pre-trained models, or be training your own model?  }

\wss{The implementation decisions can, and likely will, change over the course
of the project.  The initial documentation should be written in an abstract way;
it should be agnostic of the implementation choices, unless the implementation
choices are project constraints.  However, recording our initial thoughts on
implementation helps understand the challenge level and feasibility of a
project.  It may also help with early identification of areas where project
members will need to augment their training.}

Topics to discuss include the following:

\begin{itemize}
\item Specific programming language
\item Specific libraries
\item Pre-trained models
\item Specific linter tool (if appropriate)
\item Specific unit testing framework
\item Investigation of code coverage measuring tools
\item Specific plans for Continuous Integration (CI), or an explanation that CI
  is not being done
\item Specific performance measuring tools (like Valgrind), if
  appropriate
\item Tools you will likely be using?
\end{itemize}

\wss{git, GitHub and GitHub projects should be part of your technology.}

\section{Use of AI and LLMs}

\wss{This section is a continuation of the previous section.  The role of LLMs
in your project as important enough to warrant their own section.  
How will your team be using LLMs for code and documentation development?  
What tools will you be using?  How will you verify the output of your LLMs?}

\section{Coding Standard}

\wss{What coding standard will you adopt?}

\newpage{}

\section{Appendix --- Generative AI Use Disclosure}

% \wss{ChatGPT (version 5) was used to brainstorm ideas for the template for this disclose section.}

% \wss{This section is about the use of AI to write this document, not your
% planned use of AI for your project.  You discuss that topic in one of the
% sections in the body of this report.}

\subsection{AI Tools Used}

% \wss{Give the name of the AI tool(s) used and the version.}

\begin{itemize}
    \item ChatGPT (GPT-5.6 Luna)
\end{itemize}

\subsection{How and Where Used}

% \wss{Explain what the tool(s) were used for.  Examples include: brainstorming,
% explaining a technical concept, reviewing the document, generating or modifying
% text, generating or modifying diagrams, identifying errors or omissions or
% inconsistencies.}

% \wss{For each of the uses, identify which parts of the document were affected.}

% \wss{You do not need to report every interaction, but you should disclose the 
% uses that materially contributed to your work}

\begin{itemize}
  \item In sections 1--4 and section 16 (Team Charter), the ideas and content
  were developed by the team. GPT-5.6 Luna was used to improve wording, grammar, and flow.
\end{itemize}

\subsection{Verification of AI-Generated Content}

% \wss{\begin{itemize}
%     \item Describe how the team checked AI-generated or AI-assisted material
%     before including it in the document.
%     \item In particular, verify factual claims, technical assertions,
%     requirements, calculations, references, and other information for which an
%     AI system may produce incorrect or fabricated results.
%     \item \textbf{The team is responsible for the accuracy, completeness,
% consistency, and appropriateness of the submitted document, regardless of
% whether material was generated by a human or an AI system.}
% \end{itemize}}

All AI-assisted text was reviewed by the team before being included in the document.
The team checked that the revised wording accurately represented the team's decisions,
policies, and intended meaning. Any changes that did not reflect the team's decisions
were removed or revised.

\newpage{}

\section{Appendix --- Reflection}

% \wss{Not required for CAS 741}

% \input{../Reflection.text}

% \begin{enumerate}
%     \item Why is it important to create a development plan prior to starting the
%     project?
%     \item In your opinion, what are the advantages and disadvantages of using
%     CI/CD?
%     \item What disagreements did your group have in this deliverable, if any,
%     and how did you resolve them?
% \end{enumerate}

\subsection*{Mohammad Bilal}

\subsubsection*{Why is it important to create a development plan prior to starting the project?}

Creating a development plan gives the team a clear idea of how the project will be organized
before development starts. It helps establish expectations for how the team will communicate,
divide work, manage the repository, and handle deadlines. Without a plan, some of these decisions
would have to be made as problems come up, which could lead to confusion or inconsistent approaches
between team members. The development plan also gives the team something to refer back to throughout
the project and helps make sure everyone has a shared understanding of how the project should be managed.

\subsubsection*{In your opinion, what are the advantages and disadvantages of using CI/CD?}

One of the main advantages of CI/CD is that it provides an automated way to check new changes
before they are merged into the project. Running tests automatically can catch bugs early and
help make sure that a new change does not break functionality that was already working. This is
especially useful when multiple people are working on the same codebase, since changes from different
team members can be tested consistently. CI/CD can also save time in the long run by reducing the
amount of manual testing needed for common checks.

The main disadvantage is the additional setup and maintenance required. The team needs to spend time
creating the pipeline and making sure it continues to work as the project changes. Failures in the
pipeline can also take time to investigate, especially when the issue is with the CI/CD configuration
rather than the application itself.

\subsubsection*{What disagreements did your group have in this deliverable, if any, and how did you resolve them?}

There were no major disagreements during this deliverable. The main difference in opinion was
how GitHub Issues should be structured. One preference was to use higher-level issues, while
another was to break them down into more detailed tasks. The team discussed both approaches,
agreed on a middle ground, and used a vote to finalize the decision.

\subsection*{Jake Finlay}

\subsubsection*{Why is it important to create a development plan prior to starting the project?}
Without a development plan, collaboration can quickly become messy. Each team member tends to
bring their own habits for code formatting, pull requests, and commit messages, and nothing ends
up being consistent. A standardized workflow makes it much easier to find information in the
repository and to contribute in a meaningful way. The plan also acts as a single source of truth
for everything related to development. When disagreements come up, the team can point back to
what was already agreed upon instead of debating it again. This keeps everyone aligned for the
duration of the project.

\subsubsection*{In your opinion, what are the advantages and disadvantages of using CI/CD?}
My favourite advantage of CI/CD is automation. Simply pushing to a branch can trigger PDF
generation, builds, deployments, tests, formatting, and linting without any manual effort. The
results of these checks are visible to everyone, so reviewers can approve pull requests with much
more confidence. However, CI/CD also has some disadvantages. A poorly
configured pipeline can take a long time to run, which becomes frustrating when checks must pass
before a merge is allowed. Pipelines can also break, and a broken check may let quality issues slip
through or block progress entirely. As a result, someone has to be responsible for maintaining the
pipeline. That maintenance draws on team resources that could otherwise go toward development.

\subsubsection*{What disagreements did your group have in this deliverable, if any, and how did you resolve them?}
Our group did not have any real disagreements during this deliverable. There was some back-and-forth
about the contents of the development plan in our team group chat, but these discussions
were resolved quickly. Going forward, I expect that clear communication and good documentation
will help us settle any development-related issues that arise.

\newpage{}

\section{Appendix --- Team Charter}

\subsection*{External Goals}

The team's main goal is to develop a project that is useful, well-designed, and
worth presenting at the end of the year. Since the project addresses a real clinical problem,
the final product should reflect the needs of its users and stakeholders rather than focusing
only on the technical implementation. The project will also provide the team with experience
working with clinicians and developing software for a real-world setting. A secondary goal
is to achieve a strong academic result in the course.

\subsection*{Attendance}

\subsubsection*{Expectations}

Team members are expected to attend scheduled meetings on time and stay for the full meeting
unless they have communicated otherwise. Members should come prepared by reviewing any material
relevant to the meeting and completing work that was assigned beforehand. If someone cannot attend,
they should notify the team as early as possible and review the meeting notes afterward.
Missing meetings occasionally is understandable, but missing three consecutive meetings without
a valid reason will result in the issue being discussed with the team.

\subsubsection*{Acceptable Excuse}

Acceptable reasons for missing a meeting or deadline include illness, medical appointments,
family emergencies, academic obligations that cannot be rescheduled, and other significant
personal circumstances. Unexpected work-related commitments may also be considered when they
cannot reasonably be avoided. Whenever possible, the team member should provide notice before
the conflict occurs.

Missing a meeting or deadline due to a lack of preparation, forgetting, or a conflict that
could have been avoided will not be considered an acceptable excuse. A member who misses a
meeting or deadline is responsible for reviewing what was missed and completing any outstanding work.

\subsubsection*{In Case of Emergency}

If an emergency prevents a member from attending a meeting or completing assigned work,
they should notify the team through the team Discord chat as soon as they are able. If the member
is unable to communicate, notification should be provided when they are able to do so. The
team will determine whether the affected work can be redistributed temporarily or whether
deadlines need to be adjusted. If the situation continues for more than two weeks and
affects the team's ability to meet project deadlines, the teaching assistants and course
instructors will be contacted to determine an appropriate approach.

\subsection*{Accountability and Teamwork} 

\subsubsection*{Quality}

Team members are expected to take responsibility for their assigned work and ask for help when needed.
Assigned tasks should be completed before the relevant team meeting so that progress can be reviewed
and discussed. Every contribution to the repository will be reviewed by at least one other team member
before it is merged. Reviews should provide meaningful feedback and ensure that the work meets the project
requirements and agreed-upon standards. Major deliverables will be reviewed by the team against the course
rubric before submission, with the goal of meeting the Level 4 criteria.

The team's internal deadline for deliverables will be two days before the official submission deadline.
This will provide sufficient time to review the deliverable, address feedback, and make any necessary
revisions before submission.

As the project involves clinical stakeholders, particular care will be taken when handling information
or feedback obtained from them. Any concerns regarding privacy, stakeholder feedback, project requirements,
or decisions that may affect the users or functionality of the application should be communicated to the team
and addressed before proceeding.

\subsubsection*{Attitude}

Team members are expected to maintain a respectful, positive, and collaborative attitude throughout the project. 
Everyone should contribute their ideas, listen to the opinions of others, and be open to constructive feedback.
Disagreements are expected to be handled respectfully and focused on the work rather than the individual.
Team members should also communicate openly when they have concerns or need assistance so that issues can
be addressed without affecting the team dynamic.

\subsubsection*{Stay on Track}

The team will stay on track through regular check-ins during team meetings to discuss completed work,
current tasks, and any issues that may affect upcoming deadlines. GitHub Issues and Projects will be
used to break down larger milestones into smaller tasks, assign responsibilities, and track progress.

The team will aim to complete deliverables two days before the official deadline. This will allow time
to review the work, compare it against the course rubric, address concerns, and make any necessary revisions.

Members are expected to complete their assigned tasks within the agreed-upon timeframe and communicate
any delays as early as possible. If a member's contribution is consistently below expectations,
the issue will first be discussed with the member. If the situation does not improve, the team
may seek advice from the teaching assistant and, if necessary, the course instructor. Failure
to contribute a fair share of the work may also be reflected in peer evaluations.

\textbf{Target Metrics:}
\begin{itemize}
  \item Attendance: 90\% of scheduled team meetings.
  \item Issues: All members must contribute to issues for each deliverable, with
  their contributions reflecting approximately 20\% of the team's total workload.
  \item Code review: Every repository contribution must be reviewed by at least one other team member before being merged.
\end{itemize}

If a member does not meet a target, they should communicate the reason to the team and identify
how the issue will be addressed. If targets are repeatedly not met, the team will meet with the
teaching assistant to discuss how to improve participation and maintain an equitable distribution of work.

\subsubsection*{Team Building}

The team will have a monthly social activity to maintain team cohesion throughout the project.
The activity will be planned based on the availability and preferences of the team members.

\subsubsection*{Decision Making}

The team will make decisions through discussion and majority vote. Each member will have the
opportunity to share their opinion before a vote is held. If a disagreement cannot be resolved
within the team, a neutral mediator, either a team member or someone outside the team such as
a teaching assistant, will be consulted.

\end{document}